由基变换 [1, 8.4]，
{\small
\begin{align*}
T_x\bigl(R^ip_{1+}\bigl(T\otimes_{\calO_Z}p_2^+(\D_\lambda)\bigr)\bigr)
&=i_x^+\bigl(R^ip_{1+}\bigl(T\otimes_{\calO_Z}p_2^+(\D_\lambda)\bigr)\bigr)\\
&=R^ir_+\bigl(i_F^+\bigl(T\otimes_{\calO_Z}p_2^+(\D_\lambda)\bigr)\bigr)
=R^ir_+\bigl(i_F^+(T)\otimes_{\calO_F}(p_2\circ i_F)^+(\D_\lambda)\bigr)\\
&=R^ir_+\bigl(T_F\otimes_{\calO_F}\D_\lambda|_F\bigr).
\end{align*}}
此外，$F$ 是 $G$ 的某个一维幂幺子群的轨道，故齐性可逆 $\calO_F$-模 $T_F$ 同构于 $\calO_F$；而作为齐性扭微分算子层，$\D_\lambda|_F$ 同构于 $\D_F$。最后，当 $i<0$ 时 $R^ir_+(\D_F)=0$。至此 3.9 的证明完成。

现在回到一般情形。下一步对单根 $\alpha\in\Pi$ 所附交织函子的分析至关重要。由于态射 $p_1$ 是仿射满淹没，可用 de Rham 复形（……）把 $I_{s_\alpha}(V)$ 作为 $\calO_X$-模与 $U(\g)$-模算出来。对 $V\in\mathcal{M}_{qc}(\D_\lambda)$，$I_{s_\alpha}(V)$ 是复形 $p_{1*}(C_{Z|X}(V))$ 的第 0 个上同调，其中 $C_{Z|X}(V)$ 记复形
$$\cdots\longrightarrow0\longrightarrow T\otimes_{\calO_Z}p_2^+(V)\longrightarrow p_2^+(V)\longrightarrow0\longrightarrow\cdots.$$
特别地，$I_{s_\alpha}(V)$ 是 $p_{1*}(p_2^*(V))$ 的商。因此有自然 $U(\g)$-模态射，从 $p_{1*}(p_2^*(V))$ 的整体截面映入 $\Gamma(X,I_{s_\alpha}(V))$。由于 $\Gamma(X,p_{1*}(p_2^*(V)))=\Gamma(Z,p_2^*(V))$，这给出自然 $U_\theta$-模态射，从 $\Gamma(X,V)$ 映入 $\Gamma(X,I_{s_\alpha}(V))$。它诱导自然 $\D_{s_\alpha\lambda}$-模态射，从 $\Delta_{s_\alpha\lambda}(\Gamma(X,V))$ 映入 $I_{s_\alpha}(V)$；也就是说，得到从函子 $\Delta_{s_\alpha\lambda}\circ\Gamma$ 到 $I_{s_\alpha}$ 的态射。把这段讨论用于特例 $V=\D_\lambda$，由 C.6.1(i) 得到自然 $\D_{s_\alpha\lambda}$-模态射，从 $\D_{s_\alpha\lambda}$ 映入 $I_{s_\alpha}(\D_\lambda)$。

对 $\alpha\in\Pi$，称 $\lambda\in\h^*$ 是 $\alpha$-反支配的，若 $\check\alpha(\lambda)$ 不是严格正整数。

\begin{lemma}\label{lem:03.3.11}
设 $\alpha\in\Pi$，$\lambda\in\h^*$ 为 $\alpha$-反支配。则从 $\D_{s_\alpha\lambda}$ 到 $I_{s_\alpha}(\D_\lambda)$ 的自然态射是 $\D_{s_\alpha\lambda}$-模同构。
\end{lemma}

\begin{proof}
断言是局部的，故可应用前面的讨论。这把问题化归到 $\g=\mathfrak{sl}(2,\C)$ 情形的相应结果。此时只有一个单根 $\alpha$，且 $s_\alpha=-1$；故可置 $I=I_{s_\alpha}$。

先断言：从 $\D_{-\lambda}$ 到 $I(\D_\lambda)$ 的自然态射不为零。为此考虑从 $\Gamma(X,\D_{-\lambda})$ 到 $\Gamma(X,I(\D_\lambda))$ 的态射。只需说明截面 $1\in U_\theta=\Gamma(X,\D_{-\lambda})$ 总映到 $\Gamma(X,I(\D_\lambda))$ 的非零截面。注意 $I(\D_\lambda)$ 是复形
$$\cdots\longrightarrow0\longrightarrow p_{1*}\bigl(T\otimes_{\calO_Z}p_2^+(\D_\lambda)\bigr)\longrightarrow p_{1*}\bigl(p_2^+(\D_\lambda)\bigr)\longrightarrow0\longrightarrow\cdots$$
的第 0 个上同调，而由 3.9，该复形的其余上同调全为零。

因 $Z$ 是仿射簇\index{fangshecu@仿射簇|hyperpage}，由 Leray 谱序列可见，这是 $I(\D_\lambda)$ 的一个左消解，各项都是 $\Gamma(X,-)$-零调的 $\calO_X$-模。因此态射
$$d:\Gamma\bigl(T\otimes_{\calO_Z}p_2^+(\D_\lambda)\bigr)\longrightarrow\Gamma\bigl(Z,p_2^+(\D_\lambda)\bigr)$$
是单射，而 $\Gamma(X,I(\D_\lambda))$ 是 $d$ 的余核。从 $U_\theta=\Gamma(X,\D_{-\lambda})$ 到 $\Gamma(X,I(\D_\lambda))$ 的态射由自然态射 $U_\theta=\Gamma(X,\D_\lambda)\to\Gamma(Z,p_2^+(\D_\lambda))$ 诱导。因此 $1\in U_\theta$ 映到 $1\in\Gamma(Z,p_2^+(\D_\lambda))$ 在商映射下的像；于是只需说明 $1\in\Gamma(Z,p_2^+(\D_\lambda))$ 不在 $d$ 的像中。

为证此，利用 $d$ 是 $G$-等变这一事实，分析 $G$ 在 $\Gamma(Z,T\otimes_{\calO_Z}p_2^+(\D_\lambda))$ 上的作用。若按次数滤过 $\D_\lambda$，就得到由 $G$-齐性 $\calO_X$-模构成的滤过 $F_p\D_\lambda$（$p\in\Z_+$），且 $\Gr\D_\lambda=S(T_X)$。于是 $T\otimes_{\calO_Z}p_2^+(\D_\lambda)$ 有由 $G$-齐性 $\calO_Z$-模构成的滤过
$$F_p\bigl(T\otimes_{\calO_Z}p_2^+(\D_\lambda)\bigr)=T\otimes_{\calO_Z}p_2^*(F_p\D_\lambda),\qquad p\in\Z_+,$$
并且
$$\Gr\bigl(T\otimes_{\calO_Z}p_2^+(\D_\lambda)\bigr)=T\otimes_{\calO_Z}p_2^*\bigl(S(T_X)\bigr)=T\otimes_{\calO_Z}S\bigl(p_2^*(T_X)\bigr).$$
对次数归纳可见 $T\otimes_{\calO_Z}p_2^+(\D_\lambda)$ 的高次上同调为零，故它的滤过诱导其整体截面的滤过，且
$$\Gr\Gamma\bigl(Z,T\otimes_{\calO_Z}p_2^+(\D_\lambda)\bigr)=\Gamma\bigl(Z,T\otimes_{\calO_Z}p_2^*(S(T_X))\bigr).$$
由于群 $G$ 约化\index{yuehua@约化|hyperpage}，代数 $G$-模 $\Gamma(Z,T\otimes_{\calO_Z}p_2^+(\D_\lambda))$ 与 $\Gamma(Z,T\otimes_{\calO_Z}S(p_2^*(T_X)))$ 同构。另一方面，$T_X=\calO(-\alpha)$，故由 3.3(i)，$S(p_2^*(T_X))$ 是诸 $p_1^*(\calO(k\alpha))$（$k\in\Z_+$）的直和。此外 $T=p_1^*(\calO(\alpha))$，最终得到
$$\Gamma\bigl(Z,T\otimes_{\calO_Z}p_2^+(\D_\lambda)\bigr)=\bigoplus_{k=1}^{\infty}\Gamma\bigl(Z,p_1^*(\calO(k\alpha))\bigr).$$
\begin{annote}
底本此处印作 $T_X=\calO(-\alpha)$，但紧接着就说 $S(p_2^*(T_X))$ 是诸 $p_1^*(\calO(k\alpha))$（$k\in\Z_+$）的直和，并最终写出 $\bigoplus_{k\ge1}\Gamma(Z,p_1^*(\calO(k\alpha)))$，这要求 $T_X=\calO(\alpha)$；而在 $\mathbb{P}^1$ 上切层本也是 $T_X=\calO(2\rho)=\calO(\alpha)$（因 $2\rho=\alpha$）。$\calO(-\alpha)$ 应为余切层，而非切层。此处照录原文，未作改动，存疑。[来源：底本 p.54；参见本节 3.3 之前对 $T_w=p_1^*(\calO(\rho-w\rho))$ 的约定]
\end{annote}
由 Frobenius 互反律，群 $G$ 在 $G$-模 $\Gamma(Z,T\otimes_{\calO_Z}p_2^+(\D_\lambda))$ 的任何子模上都不平凡地作用。这说明 $1$ 不在 $d$ 的像中。

下面说明：当 $\lambda$ 为 $\alpha$-反支配时，从 $\D_{-\lambda}$ 到 $I(\D_\lambda)$ 的自然态射是同构。首先注意，$\D_{-\lambda}$ 与 $I(\D_\lambda)$ 都是 $G$-齐性 $\calO_X$-模，且自然态射 $G$-等变。故只需说明该态射在 $X$ 的基点 $x_0$ 处诱导几何纤维的同构。$\D_{-\lambda}$ 的元素右乘，在 $T_{x_0}(\D_{-\lambda})$ 上诱导左 $U(\g)$-模结构，同构于 $U(\g)\otimes_{U(\bb_{x_0})}\C_{\lambda-\rho}$。因此 $T_{x_0}(\D_{-\lambda})$ 是诸一维权空间（对应权 $\{\lambda-\rho-k\alpha\mid k\in\Z_+\}$）的直和；若 $\alpha$-反支配条件成立，它不可约。此外，$\Gamma(X,\D_\lambda)$ 的元素右乘，在 $T_{x_0}(I(\D_\lambda))$ 上诱导左 $U(\g)$-模结构，而几何纤维之间的映射是 $U(\g)$-模态射。故若 $\lambda$ 为 $\alpha$-反支配，从 $T_{x_0}(\D_{-\lambda})$ 到 $T_{x_0}(I(\D_\lambda))$ 的态射是单射。

3.9 的证明中的论证还说明
$$T_{x_0}\bigl(I(\D_\lambda)\bigr)=R^0r_+(\D_F)=\Gamma(F,\calO_F)$$
作为向量空间。另一方面，它有自然的 $U(\bb_{x_0})$-模结构，由线性形式 $-\lambda+\rho$ 给出。此外，$x_0\in X$ 的稳定化子 $B_{x_0}$ 作用在这个模上，并在 $\Gamma(F,\calO_F)$ 上诱导自然作用。因此它是诸一维权空间（对应权 $\{-k\alpha\mid k\in\Z_+\}$）的直和。$U(\bb_{x_0})$ 的元素在 $\D_\lambda|_F$ 上右乘，在 $T_{x_0}(I(\D_\lambda))$ 上诱导左 $U(\bb_{x_0})$-作用，它是第二个作用与第一个作用之差。故 $T_{x_0}(I(\D_\lambda))$ 是诸一维权空间（对应权 $\{\lambda-\rho-k\alpha\mid k\in\Z_+\}$）的直和。由此可见几何纤维之间的态射是同构。
\end{proof}

对任意 $S\subset\Sigma^+$，称 $\lambda\in\h^*$ 是 $S$-反支配的，若它对一切 $\alpha\in S$ 都是 $\alpha$-反支配的。对任意 $w\in W$，置
$$\Sigma^+_w=\{\alpha\in\Sigma^+\mid w\alpha\in-\Sigma^+\}=\Sigma^+\cap\bigl(-w^{-1}(\Sigma^+)\bigr).$$

\begin{lemma}\label{lem:03.3.12}
\begin{enumerate}[label=(\roman*),leftmargin=2.2em,itemsep=0pt,topsep=2pt]
\item $\Sigma^+_{w^{-1}}=-w(\Sigma^+_w)$。
\item 设 $w,w'\in W$ 满足 $\ell(w'w)=\ell(w')+\ell(w)$。则 $\Sigma^+_{w'w}=w^{-1}(\Sigma^+_{w'})\cup\Sigma^+_w$。
\item 设 $w,w'\in W$ 满足 $\ell(w'w)=\ell(w')+\ell(w)$。若 $\lambda\in\h^*$ 是 $\Sigma^+_{w'w}$-反支配的，则 $w\lambda$ 是 $\Sigma^+_{w'}$-反支配的。
\end{enumerate}
\end{lemma}

\begin{proof}
(i) 直接由 $\Sigma^+_w$ 的定义得出。(ii) 由 [5, Ch. VI, §1, no. 6, Cor. 2. of Prop. 17] 得出。(iii) 立即由 (ii) 得出。
\end{proof}

\begin{lemma}\label{lem:03.3.13}
设 $\lambda\in\h^*$ 反支配。则对任意 $w\in W$，$J_w\bigl(\D(\D_\lambda)\bigr)=\D(\D_{w\lambda})$。
\end{lemma}

\begin{proof}
由 3.11，命题对单反射成立。对 $\ell(w)$ 归纳证一般情形。设 $\ell(w)=k+1$。则有 $\alpha\in\Pi$ 与 $w'\in W$，使 $w=s_\alpha w'$ 且 $\ell(w')=k$。由 3.12(iii)，$w'\lambda$ 是 $\alpha$-反支配的，故由 3.7、3.9、3.11 与归纳假设，
$$J_w\bigl(\D(\D_\lambda)\bigr)=J_{s_\alpha}\bigl(J_{w'}(\D(\D_\lambda))\bigr)=J_{s_\alpha}\bigl(\D(\D_{w'\lambda})\bigr)=\D(\D_{w\lambda}).$$
\end{proof}

\begin{theorem}\label{thm:03.3.14}
设 $w\in W$，$\lambda\in\h^*$。则：
\begin{enumerate}[label=(\roman*),leftmargin=2.2em,itemsep=0pt,topsep=2pt]
\item 对任意 $V^\bullet\in\D^b(\D_\lambda)$，有 $LI_w(V^\bullet)=J_w(V^\bullet)$；
\item $I_w$ 的左上同调维数不超过 $\ell(w)$。
\end{enumerate}
\end{theorem}

\begin{proof}
显然 (ii) 由 (i) 得出。为证 (i)，须说明：对任意投射对象 $P\in\mathcal{M}_{qc}(\D_\lambda)$，当 $i\neq0$ 时 $H^i(J_w(\D(P)))=0$。首先注意，由 3.5，只需考虑 $\lambda$ 正则反支配的情形。此时由 1.3，$P$ 是某个投射 $U_\theta$-模的局部化，从而是 $(\D_\lambda)^{(J)}$ 的直和项。故只需处理 $P=\D_\lambda$ 的情形，此时可应用 3.13。
\end{proof}

作为立即可得的推论，交织函子 $LI_w$（$w\in W$）可延拓为从 $\D(\D_\lambda)$ 到 $\D(\D_{w\lambda})$ 的函子。下面证明交织函子乘积公式的一个预备版本。

\begin{lemma}\label{lem:03.3.15}
设 $w,w'\in W$ 满足 $\ell(w'w)=\ell(w')+\ell(w)$。则对任意 $\lambda\in\h^*$，从 $\D^-(\D_\lambda)$ 到 $\D^-(\D_{w'w\lambda})$ 的函子 $LI_{w'}\circ LI_w$ 与 $LI_{w'w}$ 同构。
\end{lemma}

\begin{proof}
由 3.8，只需说明：对任意投射对象 $P\in\mathcal{M}_{qc}(\D_\lambda)$，$\D_{w\lambda}$-模 $I_w(P)$ 是 $I_{w'}$-零调的。由 3.5，还可设 $\lambda$ 正则反支配。此时与 3.14 的证明一样，可设 $P=\D_\lambda$。于是由 3.13，$I_w(\D_\lambda)=\D_{w\lambda}$，断言由反复应用 3.13 与 3.14 得出。
\end{proof}

最后得到下述结果，它阐明交织函子的作用。

\begin{theorem}\label{thm:03.3.16}
设 $w\in W$，$\lambda\in\h^*$ 是 $\Sigma^+_w$-反支配的。则从 $\D^-(U_\theta)$ 到 $\D^-(\D_{w\lambda})$ 的函子 $LI_w\circ L\Delta_\lambda$ 与 $L\Delta_{w\lambda}$ 同构。
\end{theorem}

\begin{proof}
对 $\ell(w)$ 归纳。先设 $w$ 是 $\alpha\in\Pi$ 对应的单反射。任一 $V^\bullet\in\D^-(U_\theta)$ 都同构于由自由 $U_\theta$-模组成的复形 $F^\bullet$。此外 $L\Delta_\lambda(V^\bullet)=\Delta_\lambda(F^\bullet)$，而由 3.9、3.11、3.14 与 C.6.1(i)，前面描述的自然态射是从
$$L\Delta_{s_\alpha\lambda}(V^\bullet)=\Delta_{s_\alpha\lambda}(F^\bullet)=\Delta_{s_\alpha\lambda}\bigl(\Gamma(X,\Delta_\lambda(F^\bullet))\bigr)$$
到
$$I_{s_\alpha}\bigl(\Delta_\lambda(F^\bullet)\bigr)=LI_{s_\alpha}\bigl(\Delta_\lambda(F^\bullet)\bigr)=LI_{s_\alpha}\bigl(L\Delta_\lambda(V^\bullet)\bigr)$$
的同构。

设命题对一切满足 $\ell(w')\le k$ 的 $w'$ 成立。取 $w\in W$ 使 $\ell(w)=k+1$。则有 $\alpha\in\Pi$ 与 $w'\in W$，使 $w=s_\alpha w'$ 且 $\ell(w')=k$。由 3.12(iii)，$w'\lambda$ 是 $\alpha$-反支配的。此外，由归纳假设与 3.15 可见
$$LI_w\circ L\Delta_\lambda=(LI_{s_\alpha}\circ LI_{w'})\circ L\Delta_\lambda=LI_{s_\alpha}\circ(LI_{w'}\circ L\Delta_\lambda)$$
同构于 $LI_{s_\alpha}\circ L\Delta_{w'\lambda}$。于是断言由单反射情形得出。
\end{proof}

\begin{corollary}\label{cor:03.3.17}
设 $\lambda\in\h^*$ 是 $\Sigma^+_w$-反支配的，$F\in\mathcal{M}(U_\theta)$ 是平坦 $U_\theta$-模。则局部化 $\Delta_\lambda(F)$ 是 $I_w$-零调的 $\Delta_\lambda$-模。
\end{corollary}

\begin{proof}
因 $F$ 是平坦模，其高次局部化为零。故断言由与 3.18 相联系的谱序列得出。
\end{proof}

下面结果给出交织函子乘积公式的最终形式。

\begin{theorem}\label{thm:03.3.18}
设 $w,w'\in W$ 满足 $\ell(w'w)=\ell(w')+\ell(w)$。则对任意 $\lambda\in\h^*$，从 $\D(\D_\lambda)$ 到 $\D(\D_{w'w\lambda})$ 的函子 $LI_{w'}\circ LI_w$ 与 $LI_{w'w}$ 同构。
\end{theorem}

\begin{proof}
由 3.5，可设 $\lambda$ 反支配且正则。此时由 1.4 与 1.12，任一复形 $V^\bullet\in\D(\D_\lambda)$ 都拟同构于复形 $P^\bullet\in\D(U_\theta)$（由投射 $U_\theta$-模组成）的局部化 $\Delta_\lambda(P^\bullet)$。因此由 3.12、3.16 与 3.17，
$$LI_w(V^\bullet)=LI_w\bigl(\Delta_\lambda(P^\bullet)\bigr)=\Delta_{w\lambda}(P^\bullet),$$
而 $\Delta_{w\lambda}(P^\bullet)$ 由 $I_{w'}$-零调的 $\D_{w\lambda}$-模组成。于是
\begin{align*}
LI_{w'}\bigl(LI_w(V^\bullet)\bigr)&=LI_{w'}\bigl(\Delta_{w\lambda}(P^\bullet)\bigr)=I_{w'}\bigl(\Delta_{w\lambda}(P^\bullet)\bigr)\\
&=\Delta_{w'w\lambda}(P^\bullet)=LI_{w'w}\bigl(\Delta_\lambda(P^\bullet)\bigr)=LI_{w'w}(V^\bullet),
\end{align*}
定理得证。
\end{proof}

此外，若再设 $\lambda$ 正则，则有：

\begin{theorem}\label{thm:03.3.19}
设 $w\in W$，$\lambda\in\h^*$ 是 $\Sigma^+_w$-反支配的且正则。则
\begin{enumerate}[label=(\roman*),leftmargin=2.2em,itemsep=0pt,topsep=2pt]
\item $LI_w$ 是范畴 $\D(\D_\lambda)$ 与 $\D(\D_{w\lambda})$ 的等价，且同构于 $L\Delta_{w\lambda}\circ R\Gamma$；
\item 从 $\D(U_\theta)$ 到 $\D(\D_{w\lambda})$ 的函子 $LI_w\circ L\Delta_\lambda$ 与 $L\Delta_{w\lambda}$ 同构。
\end{enumerate}
\end{theorem}

\begin{proof}
先证 (ii)。任一复形 $V^\bullet\in\D(U_\theta)$ 都拟同构于由投射 $U_\theta$-模组成的复形 $P^\bullet$。因此 $L\Delta_\lambda(V^\bullet)=\Delta_\lambda(P^\bullet)$，且由 3.16 与 3.17，
$$LI_w\bigl(L\Delta_\lambda(V^\bullet)\bigr)=I_w\bigl(\Delta_\lambda(P^\bullet)\bigr)=\Delta_{w\lambda}(P^\bullet)=L\Delta_{w\lambda}(V^\bullet).$$
(i) 由 (ii) 与 1.12 得出。
\end{proof}

\begin{theorem}\label{thm:03.3.20}
设 $w\in W$，$\lambda\in\h^*$。则：
\begin{enumerate}[label=(\roman*),leftmargin=2.2em,itemsep=0pt,topsep=2pt]
\item $LI_w$ 是范畴 $\D(\D_\lambda)$ 与 $\D(\D_{w\lambda})$ 的等价；
\item $LI_w$ 是范畴 $\D^b(\D_\lambda)$ 与 $\D^b(\D_{w\lambda})$ 的等价；
\item $LI_w$ 是范畴 $\D^b_{coh}(\D_\lambda)$ 与 $\D^b_{coh}(\D_{w\lambda})$ 的等价。
\end{enumerate}
\end{theorem}

\begin{proof}
先设 $\lambda$ 正则反支配。此时由 1.12 与 3.19(i)，函子 $LI_w$ 等价于 $L\Delta_{w\lambda}\circ R\Gamma$，其逆函子等价于 $L\Delta_\lambda\circ R\Gamma$。这证明了 (i)。因函子 $\Gamma$ 的右上同调维数有限，而对正则 $\mu$，由 1.13，局部化函子 $\Delta_\mu$ 的左上同调维数有限，(ii) 得证。最后一条由 1.16 得出。

一般情形由 3.5 用几何平移得出。
\end{proof}

下面改进交织函子左上同调维数的估计。

\begin{theorem}\label{thm:03.3.21}
设 $w\in W$，$\lambda\in\h^*$。则 $I_w:\mathcal{M}_{qc}(\D_\lambda)\to\mathcal{M}_{qc}(\D_{w\lambda})$ 的左上同调维数不超过 $\operatorname{Card}(\Sigma^+_w\cap\Sigma_\lambda)$。
\end{theorem}

\begin{proof}
由 3.5，可设 $\lambda$ 正则反支配。对 $\ell(w)$ 归纳。先设 $w$ 是 $\alpha\in\Pi$ 对应的反射。则 $I_{s_\alpha}$ 的左上同调维数不超过 1。再设 $\check\alpha(\lambda)\notin\Z$。则 $s_\alpha(\lambda)$ 也正则反支配。由 3.19(i)，$LI_{s_\alpha}$ 是范畴 $\D(\D_\lambda)$ 与 $\D(\D_{s_\alpha\lambda})$ 的等价，且同构于 $L\Delta_{s_\alpha\lambda}\circ R\Gamma$。因由 1.3，函子 $\Gamma:\mathcal{M}_{qc}(\D_\lambda)\to\mathcal{M}(U_\theta)$ 与 $\Delta_{s_\alpha\lambda}:\mathcal{M}(U_\theta)\to\mathcal{M}_{qc}(\D_{s_\alpha\lambda})$ 都是范畴等价，故 $I_{s_\alpha}:\mathcal{M}_{qc}(\D_\lambda)\to\mathcal{M}_{qc}(\D_{s_\alpha\lambda})$ 也是范畴等价。于是其左上同调维数为 0。这完成了单反射情形。

设命题对一切满足 $\ell(w')\le k$ 的 $w'$ 成立。取 $w\in W$ 使 $\ell(w)=k+1$。则有 $\alpha\in\Pi$ 与 $w'\in W$，使 $\ell(w')=k$。由 3.12(ii)，
$$\Sigma^+_{s_\alpha w'}\cap\Sigma_\lambda=\bigl(\{w'^{-1}(\alpha)\}\cup\Sigma^+_{w'}\bigr)\cap\Sigma_\lambda,$$
故
$$\operatorname{Card}\bigl(\Sigma^+_w\cap\Sigma_\lambda\bigr)=\operatorname{Card}\bigl(\Sigma^+_{w'}\cap\Sigma_\lambda\bigr)+\operatorname{Card}\bigl(\{\alpha\}\cap\Sigma_{w'\lambda}\bigr).$$
断言由 3.18、单反射情形与归纳假设得出。
\end{proof}

\begin{corollary}\label{cor:03.3.22}
设 $w\in W$，$\lambda\in\h^*$ 满足 $\Sigma^+_w\cap\Sigma_\lambda=\emptyset$。则 $I_w:\mathcal{M}_{qc}(\D_\lambda)\to\mathcal{M}_{qc}(\D_{w\lambda})$ 是范畴等价，其逆为 $I_{w^{-1}}$。
\end{corollary}

\begin{proof}
由 3.12(i) 与 3.21，$I_w$ 与 $I_{w^{-1}}$ 都正合。又由 3.5，可设 $\lambda$ 正则反支配。由此可见 $w\lambda$ 正则且反支配。由 3.19(i)，$LI_w$ 是范畴 $\D(\D_\lambda)$ 与 $\D(\D_{w\lambda})$ 的等价，且同构于 $L\Delta_{w\lambda}\circ R\Gamma$。由于由 1.3，函子 $\Gamma:\mathcal{M}_{qc}(\D_\lambda)\to\mathcal{M}(U_\theta)$ 与 $\Delta_{w\lambda}:\mathcal{M}(U_\theta)\to\mathcal{M}_{qc}(\D_{w\lambda})$ 都是范畴等价，$I_w:\mathcal{M}_{qc}(\D_\lambda)\to\mathcal{M}_{qc}(\D_{w\lambda})$ 也是范畴等价。同样的论证适用于 $I_{w^{-1}}$。这也说明它们的复合同构于恒等函子。
\end{proof}

\begin{theorem}\label{thm:03.3.23}
设 $w\in W$，$\lambda\in\h^*$ 是 $\Sigma^+_w$-反支配的。则从 $\D(\D_\lambda)$ 到 $\D(U_\theta)$ 的函子 $R\Gamma\circ LI_w$ 与 $R\Gamma$ 同构。
\end{theorem}

\begin{proof}
若 $\lambda$ 正则，这由 1.12 得出。对奇异的 $\lambda$，可取 $w'\in W$ 与 $\nu\in P(\Sigma)$，使 $w'\lambda$ 与 $-w'\nu$ 反支配，$w'(\lambda-\nu)$ 反支配且正则，并且 $\lambda-\nu$ 是 $\Sigma^+_w$-反支配的。设 $V^\bullet\in\D(\D_\lambda)$。因 $I_w$ 的左上同调维数有限，$V^\bullet$ 拟同构于由 $I_w$-零调 $\D_\lambda$-模组成的复形 $C^\bullet$。此外由 3.5，复形 $C^\bullet(-\nu)$ 由 $I_w$-零调的 $\D_{\lambda-\nu}$-模组成，且
$$LI_w\bigl(V^\bullet(-\nu)\bigr)=I_w\bigl(C^\bullet(-\nu)\bigr)=I_w(C^\bullet)(-w\nu)=LI_w(V^\bullet)(-w\nu).$$
此外，
$$R\Gamma\bigl(LI_w(V^\bullet)(-w\nu)\bigr)=R\Gamma\bigl(LI_w(V^\bullet(-\nu))\bigr)=R\Gamma\bigl(V^\bullet(-\nu)\bigr),$$
这里用了 $\lambda-\nu$ 正则情形的结论。另一方面，由 C.2.1，若记 $F$ 为以 $w'\nu$ 为最高权的不可约有限维 $\g$-模，则有
{\small
\begin{align*}
R\Gamma\bigl(LI_w(V^\bullet)\bigr)
&=R\Gamma\Bigl(\bigl(LI_w(V^\bullet)(-w\nu)\otimes_{\calO_X}F\bigr)[\lambda]\Bigr)\\
&=\Bigl(R\Gamma\bigl(LI_w(V^\bullet)(-w\nu)\bigr)\otimes_\C F\Bigr)[\lambda]
=\Bigl(R\Gamma\bigl(V^\bullet(-\nu)\bigr)\otimes_\C F\Bigr)[\lambda]\\
&=R\Gamma\Bigl(\bigl(V^\bullet(-\nu)\otimes_{\calO_X}F\bigr)[\lambda]\Bigr)=R\Gamma(V^\bullet).
\end{align*}}
\end{proof}

最后指出下述事实。它说明，一般情形下 3.21 对交织函子左上同调维数的估计是最佳的。

\begin{proposition}\label{prop:03.3.24}
设 $w\in W$，$\lambda\in P(\Sigma)$。则
$$L^iI_w\bigl(\calO(\lambda+\rho)\bigr)=0,\qquad i\neq-\ell(w),$$
并且
$$L^{-\ell(w)}I_w\bigl(\calO(\lambda+\rho)\bigr)=\calO(w\lambda+\rho).$$
\end{proposition}

\begin{proof}
由 3.5，可设 $\lambda$ 反支配且正则。此时断言由 3.16 与 1.22 立即得出。
\end{proof}
